\begin{textsample}{POS Dim 4 – persona\_gpt – Score 28.00 – t489\_gpt.txt}  \label{ex:f4_pos_025}
Behind the \textbf{seafood} on our plates lies a reality that is too often hidden from view.
Across the global \textbf{seafood} industry, serious human rights abuses, environmental destruction, and climate impacts are deeply intertwined.
In many industrial \textbf{fishing} \textbf{operations}, especially those far from \textbf{shore}, migrant workers from Southeast Asia are among the most vulnerable.
They can face exploitation, physical and verbal abuse, and inhumane living and working \textbf{conditions}.
Long periods at \textbf{sea}, cramped quarters, lack of medical care, and dangerous work are common features of this system.
Transshipment at \textbf{sea} – where \textbf{catch} is transferred from one \textbf{vessel} to another without returning to \textbf{port} – can trap workers in isolation for months or even years, cutting them off from oversight, legal recourse, or even contact with their families.
This isolation makes it easier for abuses to remain unseen and unreported.
At the same time, the \textbf{oceans} that sustain these \textbf{fisheries} are under mounting pressure from climate change.
Rising \textbf{temperatures} and \textbf{ocean} acidification threaten marine \textbf{ecosystems}, undermining the health of \textbf{fish} populations and the communities that depend on them.
Yet healthy \textbf{oceans} and marine \textbf{life} are also powerful allies against climate change.
Marine \textbf{ecosystems} can absorb and \textbf{store} significant amounts of carbon, helping to buffer the planet from even more severe climate disruption.
Protecting \textbf{ocean} health is therefore critical not only for biodiversity and food security, but also for stabilising the climate.
Overfishing is eroding that resilience.
Decades of intense \textbf{fishing} pressure have pushed many \textbf{fish} populations to the \textbf{brink}.
Illegal, unreported and unregulated ( IUU ) \textbf{fishing} is a major driver of this crisis, with an estimated value of up to \$23.5 billion every year.
These shadowy \textbf{operations} often target vulnerable \textbf{stocks}, ignore \textbf{catch} limits, and operate with little to no regard for labour rights or environmental rules.
The result is a double blow : depleted \textbf{fish} populations and ongoing human rights risks.
The \textbf{methods} used in industrial \textbf{fishing} can be highly destructive. \textbf{Bottom} trawling drags heavy \textbf{nets} across the \textbf{seabed}, damaging fragile \textbf{habitats} and killing a wide range of marine \textbf{life}, not just the target \textbf{species}.
Longlining sets out vast lines of baited hooks that can entangle or kill non-target \textbf{species}, including seabirds, \textbf{turtles}, and \textbf{sharks}.
Purse seining with \textbf{fish} aggregating devices ( FADs ) concentrates \textbf{fish} in one place, making them easier to \textbf{catch} in huge quantities, but also increasing bycatch and putting additional pressure on already stressed populations.
These practices undermine marine biodiversity and the long-term sustainability of \textbf{fisheries}.
For coastal communities, the consequences are profound.
Many people rely on the \textbf{sea} for food, income, and cultural identity.
As \textbf{fish} \textbf{stocks} \textbf{decline} and destructive \textbf{fishing} practices spread, small-scale \textbf{fishers} often find themselves competing with industrial \textbf{fleets} for dwindling resources. \textbf{Livelihoods} become more precarious, local economies suffer, and traditional ways of \textbf{life} are put at risk.
Consumers and citizens have an important role to play in changing this system.
Choosing \textbf{seafood} from sustainable sources helps support \textbf{fisheries} that respect ecological limits and reduce harm to marine \textbf{ecosystems}.
But individual choices alone are not enough.
Public pressure on governments is essential to secure strong, enforceable rules for ethical and sustainable \textbf{fisheries} management.
That includes cracking down on IUU \textbf{fishing}, ending destructive \textbf{fishing} \textbf{methods}, and protecting the rights and safety of workers at \textbf{sea}.
By demanding transparency, accountability, and sustainability from the \textbf{seafood} industry and from those who regulate it, we can help protect both people and the \textbf{oceans} we all depend on.

% matched lemmas: bottom, brink, catch, condition, decline, ecosystem, fish, fisher, fishery, fishing, fleet, habitat, life, livelihood, method, net, ocean, operation, port, sea, seabed, seafood, shark, shore, specie, species, stock, store, temperature, turtle, vessel
\end{textsample}
